\section{恒定磁场}

\subsection{磁场}

运动电荷的周围既有电场又有磁场。磁场再宏观上表现出
\begin{enumerate}
	\item 对运动电荷有力的作用；
	\item 载流导线再磁场中移动时，磁场会对载流导线做功。
\end{enumerate}

\subsubsection{磁感应强度}

\begin{theorem}[毕奥--萨伐尔定律]
	对于一个电流元，其到某点的位矢为 $\vect{r}$，那么它对该点的磁感应强度贡献为
	$$
	\dd{\vect{B}} = \frac{\mu_0}{4\pi} \frac{I \dd{\vect{l}} \cp \vect{r}}{r^3}
	$$
\end{theorem}

而对于一个闭合载流回路 $L$，其产生的磁感应强度为一个路径积分：
$$
\vect{B} = \frac{\mu_0}{4\pi} \oint_{L} \frac{I \dd{\vect{l}} \cp \vect{r}}{r^3}
$$

\subsubsection{磁偶极子}

\begin{definition}[面载流线圈的磁矩]
	定义平面载流线圈的\textbf{磁矩}为
	$$
	\vect{p}_m = I \vect{S}
	$$
	其中 $\vect{S}$ 是右旋的方向向量。
\end{definition}

\subsection{磁场的高斯定理}

利用矢量场的高斯定理：
$$
\Phi_B = \oiint_{S} \vect{B} \vdot \dd{\vect{S}} = \vect{0},\quad \iiint_V \div \vect{B} \dd{V} = \vect{0}
$$
得到 $\div \vect{B} = \vect{0}$。

说明，恒定磁场的散度为 $0$，它是一个无源场。

\subsection{安培环路定理}

安培环路定理表述为，对于任意闭合环路 $L$，都有
$$
\oint_L \vect{B} \vdot \dd{\vect{l}} = \mu_0 I_{enc}
$$
其中 $I_{enc}$ 是穿过环路的净电流。其中净点流的正方向由右手法则决定。
